% ============================================================
%  SECTION 4 — METHODOLOGY
% ============================================================
\section{Methodology}
\label{sec:methodology}

% ------------------------------------------------------------
\subsection{Materials}
\label{subsec:materials}
% ------------------------------------------------------------

The experimental evaluation is conducted on the UCSB Breast
Cancer Histology dataset, provided by the Center for
Bio-Image Informatics of the University of California, Santa
Barbara~\cite{Gelasca2008}. The dataset comprises 58
hematoxylin and eosin~(H\&E) stained histological images of
breast tissue at a fixed resolution of $896 \times 768$
pixels, partitioned into two classes: 32 benign and 26
malignant samples. The class imbalance ratio of
approximately 1.23:1 is mild enough to not require
resampling for the fractal dimension estimation task, since
$D$ is computed independently per image and no
classifier is trained on the extracted features.

This dataset was selected for three reasons directly relevant
to the algorithmic focus of this work. First, its fixed
resolution ensures that the pixel count $N = 896 \times 768
= 688{,}128$ is identical across all images, making
wall-clock time comparisons across methods and
formulations free of resolution confounds. Second, the two
well-defined tissue classes — benign and malignant — provide
a concrete discriminatory benchmark against which the
effectiveness of each algorithm's $D$ estimate can be
assessed, following the methodology established in prior
fractal-based histological classification
work~\cite{Roberto2021,Borgue2025}. Third, the dataset is
publicly available as a Kaggle dataset
(\texttt{ucsb-breast-cancer-histology}), enabling direct
reproducibility of all experiments within the execution
environment described in Section~\ref{subsec:environment}.

All images are loaded as 8-bit grayscale arrays ($G = 256$)
prior to processing. For the binary-domain methods~(BC and
GB), a binarisation step is applied using Otsu's global
threshold~\cite{Otsu1979} before the fractal dimension
estimation loop; this step is timed and reported separately
from the algorithmic core, as described in
Section~\ref{subsec:environment}.

% ------------------------------------------------------------
\subsection{Experimental Environment and Protocol}
\label{subsec:environment}
% ------------------------------------------------------------

All experiments are executed on Kaggle Notebooks
(\url{https://www.kaggle.com}) using a CPU-only session,
to ensure that the measured execution times reflect the
sequential computational complexity of each algorithm
without GPU acceleration. The hardware configuration
provided by the Kaggle CPU session consists of an Intel
Xeon processor with two virtual cores and approximately
13~GB of available RAM. The software environment is fixed
at Python~3.10 with NumPy~1.26, SciPy~1.11, and
scikit-image~0.22; the complete \texttt{requirements.txt} is included in the
\textit{github} repository to ensure reproducibility.
\footnote{Source code available at:
\href{https://https://github.com/gabrielmbrum} {\texttt{https://github.com/gabrielmbrum}}.}

The scale set used across all four methods is fixed as
$\mathcal{S} = \{2, 4, 8, 16, 32, 64\}$, comprising $m = 6$
scales spanning approximately 1.5 decades of spatial
frequency. This range respects the lower bound imposed by
the pixel size and the upper bound of $L_{\max} = 64 \ll
\sqrt{N} \approx 830$, ensuring that the covering assumption
underlying each algorithm remains valid throughout the scale
range. The same $\mathcal{S}$ is used for all methods to
guarantee that differences in execution time and $R^{2}$
are attributable solely to algorithmic structure rather
than to differences in scale coverage.

Each algorithm — in both its naive and optimised
formulations — is executed 10 times per image. The first
run is discarded as a warm-up to eliminate cold-cache and
interpreter overhead effects; the remaining 9 runs are used
to compute the mean wall-clock time $\bar{t}$ and standard
deviation $\sigma_t$. Wall-clock time is measured using
Python's \texttt{time.perf\_counter()}, which provides
sub-microsecond resolution on the target platform. Runs
exhibiting a coefficient of variation

\begin{equation}
  CV = \frac{\sigma_t}{\bar{t}} \times 100\%
  \label{eq:cv}
\end{equation}

\noindent exceeding 5\% are flagged and re-evaluated with
additional runs until $CV < 5\%$, following the adaptive
stability criterion recommended for wall-clock
benchmarking~\cite{Flemming1986}. This threshold ensures
that the reported means are statistically stable without
imposing an unnecessarily large number of repetitions.

Pre-processing time is measured and reported separately
from the algorithmic core. Concretely, for BC and GB, the
wall-clock time of the Otsu binarisation step is recorded
independently of the subsequent coverage computation,
allowing the contribution of the pre-processing stage to
the total execution budget to be quantified and discussed
in Section~\ref{sec:results}. For DBC and BM, which
operate directly on the grayscale input, no binarisation
is performed and pre-processing time is zero by definition.

% ------------------------------------------------------------
\subsection{Evaluation Metrics}
\label{subsec:metrics}
% ------------------------------------------------------------

Three orthogonal evaluation axes are considered, each
capturing a distinct dimension of the algorithmic
trade-off under study.

\subsubsection{Execution Time}

The primary time metric is the mean wall-clock time
$\bar{t}$ (in seconds) averaged over all 58 images and
9 valid runs per image. To assess the agreement between
theoretical predictions and empirical scaling behaviour,
$\bar{t}$ is plotted as a function of image size $N$ across
sub-images of increasing resolution, and the empirical
slope of the $\log \bar{t}$ versus $\log N$ curve is
compared against the theoretical exponent derived in
Section~\ref{sec:algorithms}. Deviations between empirical
and theoretical slopes are attributed to constant factors,
memory hierarchy effects, and interpreter overhead, and are
discussed qualitatively in Section~\ref{sec:results}.

\subsubsection{Memory Consumption}

Auxiliary memory consumption is measured as the peak
heap allocation incurred by each algorithm beyond the
read-only input buffer, using Python's
\texttt{tracemalloc} module. The input buffer of size
$\Theta(N)$ bytes is excluded from the measurement by
recording the allocation differential between the start
and end of the algorithmic core. This isolates the
auxiliary space — the integral image for BC and GB, the
sparse tables for DBC, and the blanket arrays for BM —
and allows direct comparison against the $\Theta(N)$ and
$\Theta(N \log N)$ theoretical bounds derived in
Section~\ref{sec:algorithms}.

\subsubsection{Effectiveness}

Effectiveness is defined as the capacity of each
algorithm's $D$ estimate to statistically discriminate
between the benign and malignant tissue classes of the
UCSB dataset. Formally, for each method, $D$ is computed
for every image and the class-conditional distributions
$\{D_i : i \in \text{benign}\}$ and $\{D_j : j \in
\text{malignant}\}$ are obtained. The primary effectiveness
metric is the inter-class separation

\begin{equation}
  \Delta D
    = \bar{D}_{\text{malignant}} - \bar{D}_{\text{benign}},
  \label{eq:delta_d}
\end{equation}

\noindent where $\bar{D}_{\text{malignant}}$ and
$\bar{D}_{\text{benign}}$ denote the respective class
means. A larger $|\Delta D|$ indicates that the method
assigns meaningfully different fractal dimensions to the
two tissue types, reflecting the clinically established
observation that malignant tissues exhibit higher surface
roughness and structural complexity than benign
ones~\cite{Lopes2011,Rangayyan2007}. Statistical
significance of $\Delta D$ is assessed via Welch's
two-sample $t$-test at significance level $\alpha = 0.05$,
which does not assume equal variances across classes and
is appropriate given the mild class imbalance of the
dataset~\cite{Welch1947}.

Two complementary effectiveness indicators are also
reported. The coefficient of determination $R^{2}$ of the
log-log regression used to estimate $D$ from each image
quantifies the degree to which the object exhibits
consistent power-law scaling across the chosen scale set;
values close to unity are a necessary condition for
interpreting $D$ as a valid fractal descriptor~\cite{
Lopes2011}. Additionally, for the binary methods~(BC and
GB), the sensitivity of $\Delta D$ to the Otsu threshold
is evaluated by perturbing $t^{*}$ by $\pm 10\%$ and
measuring the resulting variation in $D$, providing a
quantitative characterisation of the threshold sensitivity
trade-off that distinguishes binary from grayscale methods.